\section{Analytic extension of supergravity solutions
}\label{sec:analyticext}

String theory at low energies is described by ten dimensional supergravity, or eleven dimensional supergravity in its M-theory incarnation. In suitable compactification contexts, it is possible to reduce the theory to a lower dimensional supergravity and select a subsector of fields that constitutes a consistent truncation.
This means that the solutions to the equations of motion in the consistent truncation are also solutions of the ten or eleven dimensional supergravity equations.

We focus on consistent truncations with real scalar fields that are minimally coupled to gravity and have a canonical kinetic term. Denoting the scalars by $\phi^I$ with $I=1,\dots,n$, the supergravity action of this subsector is%
\footnote{We use mostly plus signature for the metric throughout the paper. In our conventions, uppercase latin indices $M,N,\dots$ label all the directions in the bulk, while greek indices $\mu,\nu,\dots$ label the directions parallel to the boundary. When directions are only spatial, we use lowercase latin indices $i,j,\dots$.}
\begin{equation}
	S_{d+1}=\frac{1}{16\pi\, G_{d+1}}\int \dd^{d+1} x\, \sqrt{-G}\left[R-\frac{1}{2}\partial_M \phi^I\partial^M\phi^I-V\left(\{\phi^I\}\right)\right]\ ,
\end{equation}
with $G_{d+1}$ the Newton constant in $d+1$ dimensions. 

If the scalar potential is an even function of one or more scalar fields, it is possible to perform an analytic extension to imaginary values that keeps the action real. Suppose that we order the scalars so that we do the analytic extension on the first $k\leq n$ of them
\begin{equation}
\begin{split}
	&\phi^a\to \phi^a =i\, \eta^a
	\ , \qquad  
	a=1,\dots, k
	\ ,\\
	&\phi^i\to \phi^i
	\ ,\qquad\qquad\quad\,
	i=k+1,\dots,n\ .
\end{split}
\end{equation}
This analytic extension flips the sign of the kinetic term, so the scalars $\eta^a$ become ghosts. In general, this leads to violation of the energy conditions, including the null energy condition.

If the consistent truncation is supersymmetric, there exists an analytic superpotential $W\left(\{\phi^I)\}\right)$ related to the potential through
\begin{equation}
V=\frac{1}{2}\sum_{I=1}^n\left(\frac{\partial W}{\partial \phi^I}\right)^2-\frac{d}{4(d-1)}W^2\ .
\end{equation}
Supersymmetric domain wall solutions can be found using the BPS equations. The metric in domain wall coordinates is of the form
\begin{equation}
	\label{eq:domain_wall}
	\dd s^2
	=
	\dd r^2 + e^{2A(r)}\eta_{\mu\nu}\,\dd x^\mu \dd x^\nu\ ,
\end{equation}
with $\eta_{\mu\nu}$ the $d$-dimensional Minkowski metric. The scalars only depend on the $r$ coordinate. The BPS equations are
\begin{equation}
	\partial_r\phi^I=\partial_{\phi^I}W\ ,\qquad \partial_r A= -\frac{W}{2(d-1)}\ .
\end{equation}
If also the superpotential is an even function of one or several of the $\phi^I$, then the analytic extension of the BPS equations,
\begin{equation}\label{eq:BPS}
	\partial_r \eta^a=-\partial_{\eta^a}\left( W\Big|_{\phi^a=i \eta^a}\right)\ ,\quad \partial_r \phi^i =\partial_{\phi^i}W\Big|_{\phi^a=i \eta^a}\ , \quad  \partial_r A= -\frac{W}{2(d-1)}\Big|_{\phi^a=i \eta^a}\ ,
\end{equation}
gives a supersymmetric domain wall solution with real values for $\eta^a$, $\phi^i$ and $A$. If the ten or eleven dimensional uplifted metric is an even function of the analytically continued scalars, then it will also be real. The analytic extension also modifies the potential to
\begin{equation}
\label{eq:VfromW}
V=-\frac{1}{2}\sum_{a=1}^k\left(\frac{\partial W}{\partial \eta^a}\right)^2+\frac{1}{2}\sum_{i=k+1}^n\left(\frac{\partial W}{\partial \phi^i}\right)^2-\frac{d}{4(d-1)}W^2\ .
\end{equation}

A special case where the potential and superpotential are guaranteed to be even functions of a real scalar is when the scalar corresponds to the modulus of a complex field in the larger supergravity theory of which the consistent truncation is a subsector
\begin{equation}
	\Phi^a=\phi^a e^{i\theta^a}\ ,\qquad {\Phi^a}^\dagger=(\Phi^a)^*=\phi^a e^{-i\theta^a}\ .
\end{equation}
The analytic extension can be combined with a shift
\begin{equation}
	\theta^a\to \theta^a=\tilde\theta^a-\frac{\pi}{2}\ .
\end{equation}
Then, the analytic extension of the complex scalar becomes
\begin{equation}
	\Phi^a\to H^a=\eta^a e^{i\tilde\theta^a}\ ,\qquad {\Phi^a}^\dagger\to {H^a}^\dagger=-\eta^a e^{-i\tilde\theta^a}=-(H^a)^*.
\end{equation}

Another possibility is that the internal geometry has a $\mathbb{Z}_2$ symmetry, which translates into a symmetry of the action in the truncated theory under $\Phi^a\to -\Phi^a$. For instance, if the internal geometry is a fibration of two spheres of the same dimension and the scalar corresponds to the ratio of the radii
\begin{equation}
\dd s^2_{\rm int}\sim e^{\alpha}\dd\Omega_n+e^{-\alpha} \dd\Omega_n\ ,
\end{equation}
the geometry is typically symmetric under $\alpha\to -\alpha$ and the truncated theory including this field depends on even powers of it (or even combinations of exponentials).%
\footnote{
Axions usually come with a $\mathbb{Z}_2$ symmetry too, since they are pseudoscalar fields.
}

\subsection{Examples of supersymmetric domain wall solutions}
\label{subsec:superpotential}

We list here a few classic supersymmetric domain wall solutions of consistent truncations that have a superpotential which is an even function of one or several scalar fields. In Appendix~\ref{app:trunc}, a consistent truncation can be found that encompasses most of the five-dimensional models discussed in this section and explicitly shows their charged-scalar origin. 

\begin{enumerate}
\item Solutions dual to the Coulomb branch of ${\cal N}=4$ super Yang-Mills (SYM) preserving $SO(3)\times SO(3)$ symmetry. They can be found using a consistent truncation to $d+1=5$ dimensional supergravity. The superpotential and the solutions, including the uplift to type IIB supergravity, were given in \cite{Freedman:1999gk,Cvetic:1999xx}. There is a single scalar $\phi$ with a superpotential
\begin{equation}
	W(\phi)= \pm \frac{6}{L}\cosh\left(\frac{\phi}{\sqrt{3}}\right)\,.
\end{equation}
The uplift to type IIB supergravity can be interpreted as the geometry produced by a distribution of parallel D3 branes that preserves a $SO(3)\times SO(3)$ subgroup of rotations in the transverse space.

\item Solutions dual to the Coulomb branch of ABJM preserving $SO(4)\times SO(4)$ symmetry. They can be found using a consistent truncation to $d+1=4$ dimensional supergravity. The potential and solutions, including the uplift to eleven dimensional supergravity, are given in \cite{Cvetic:1999xx}. An explicit expression for the superpotential can be found in \eg \cite{Tarrio:2013qga}. There is a single scalar $\phi$ with a superpotential
\begin{equation}
	W(\phi)=\pm\frac{4}{L}\cosh \left(\frac{\phi}{2}\right)\,.
\end{equation}
The uplift to eleven dimensional supergravity can be interpreted as the geometry produced by a distribution of parallel M2 branes preserving an $SO(4)\times SO(4)$ subgroup of rotations in the transverse space.

\item Solutions dual to the ${\cal N}=2^*$ SYM theory. They can be found using a consistent truncation to $d+1=5$ dimensional supergravity. The superpotential and the solutions, including the uplift to type IIB supergravity were given in \cite{Pilch:2000ue}. The extension to a truncation with a complex scalar charged under a gauge field can be found in \cite{Bobev:2010de}. There are two scalar fields, $\chi$ (the modulus of the complex scalar) and $ \alpha$, that can be canonically normalized using $\phi^1=2\sqrt{2}\,\chi$  and $\phi^2= 2\sqrt{6}\,\alpha$. The superpotential for the canonically normalized fields is
\begin{equation}
	W(\phi^1,\phi^2)=\pm\frac{2}{L} e^{-\frac{\phi^2}{\sqrt{6}}} \left[2 + e^{\sqrt{\frac{3}{2}}\phi^2} \cosh\left(\frac{\phi^1}{\sqrt{2}}\right)\right]\ .
\end{equation}
The superpotential is even in $\phi^1$ but not in $\phi^2$, so the analytic extension to purely imaginary values is restricted to $\phi^1$. 

\item Solutions dual to the ${\cal N}=1^*$ SYM theory with a mass for one chiral multiplet (the Leigh-Strassler flow). They can be found using a consistent truncation to $d+1=5$ dimensional supergravity. The superpotential and the solutions were given in \cite{Freedman:1999gp,Pilch:2000fu}, with the latter paper also describing the uplift to type IIB supergravity. There are two scalar fields $\chi,\alpha$, that can be canonically normalized using $\phi^1=2\chi$, $\phi^2 =2\sqrt{6}\,\alpha$. The superpotential for the canonically normalized fields has the form
\begin{equation}
	W(\phi^1,\phi^2)= \pm\frac{e^{-\frac{\phi^2}{\sqrt{6}}}}{L}\left\{e^{\sqrt{\frac{3}{2}}\,\phi_2}\left[\cosh\left(\phi^1\right)-3 \right] - 2\left[1 + \cosh\left(\phi^1\right)\right]\right\}\ .
\end{equation}
Also in this case, the superpotential is even in $\phi^1$ but not in $\phi^2$. 

\item Solutions dual to the ${\cal N}=1^*$ SYM theory with equal masses for two chiral multiplets (the GPPZ solution). They can be found using a consistent truncation to $d+1=5$ dimensional supergravity. The superpotential and the solutions were given in \cite{Girardello:1999bd}. The complete uplifts to type IIB supergravity were found in \cite{Bobev:2018eer,Petrini:2018pjk}. There are two scalar fields, $m$ and $\sigma$, that can be canonically normalized using $\phi^1= 2m$, $\phi^2 = 2\sigma$. The superpotential for the canonically normalized fields is
\begin{equation}
	W(\phi^1,\phi^2) = \pm \frac{3}{L}\left[\cosh\left({\frac{\phi^1}{\sqrt{3}}}\right)+\cosh\left({\phi^2}\right)\right]\ .
\end{equation}
In this case, both scalars correspond to the modulus of complex fields in a larger truncation, so the superpotential is even in both. 

\item Solutions dual to states with spontaneously broken R-symmetry in $3+1$-dimensional ${\cal N}=1$ SCFTs \cite{Gubser:2009qm}. They can be found using a consistent truncation to $d+1=5$ dimensional supergravity. There is a single canonically normalized scalar, with a superpotential \cite{Cassani:2010uw,Gauntlett:2010vu}
\begin{equation}
	W(\phi)= \pm \frac{6}{L}\cosh^2\left(\frac{\phi}{2}\right)\ .
\end{equation}
\item Solutions dual to states with spontaneously broken R-symmetry in $2+1$-dimensional SCFTs \cite{Gauntlett:2009dn}. They can be found using a consistent truncation to $d+1=4$ dimensional supergravity. There is one scalar field $\chi$, that is related to the canonically normalized field through $\phi=2\tanh^{-1}\frac{\chi}{\sqrt{2}}$. The superpotential for the canonically normalized field is found to be
\begin{equation}
	W(\phi)= \pm \frac{2}{L}\left[\cosh\left(\phi\right)+1\right]=\pm \frac{4}{L}\cosh^2\left(\frac{\phi}{2}\right)\ .
\end{equation}

\item Solutions dual to flows triggered by vevs for boson bilinears. They correspond to the model in appendix~\ref{app:trunc} with vanishing charged scalars. The canonically normalized neutral fields are $\phi^1=2\sqrt{2}\beta$ and $\phi^2=2\sqrt{6}\alpha$. Despite the fact that $\phi^1$ is not the modulus of a complex scalar, the superpotential is an even function of it 
\begin{align}
W(\phi^1,\phi^2)=\pm\frac{2}{L}e^{-\frac{\phi^2}{\sqrt{6}}}\left[e^{\sqrt{\frac{3}{2}}\phi^2}+2\cosh\left(\frac{\phi^1}{\sqrt{2}}\right)\right]\,.
\end{align}

\end{enumerate}

\subsection{Analytic extension of domain wall solutions}
\label{subsec:examples}

The superpotentials we have listed above depend on the scalars through hyperbolic functions. When doing the analytic extension to purely imaginary values, these become trigonometric functions with some periodicity under shifts of the scalars. It is natural to identify critical points of the superpotential as equivalent if they are related by a shift by one period, but there are additional critical points also at half-periods that are unequivalent to the critical point at the origin of the target space.
The flows that are solutions to the BPS equations \eqref{eq:BPS} interpolate between the critical point at the origin of the target space of the scalars and the new critical points of the superpotential. We will use the convention for the sign of the superpotential that puts at $r\to-\infty$ the asymptotic AdS boundary corresponding to the critical point at the origin.

Here we list the solutions we find for each case. 
\begin{enumerate}
	\item The analytic extension $\phi\to \phi=i\, \eta$ makes the superpotential periodic in $\eta$, with period $2\pi\sqrt{3}$,
\begin{equation}
	\label{eq:Wex1}
	W(\eta) = \frac{6}{L} \cos\left(\frac{\eta}{\sqrt{3}}\right)\ .
\end{equation}
There is a flow between the two critical points of the superpotential at $\eta=0$ (when $r\to-\infty)$ and $\eta=\pi\sqrt{3}$ (when $r\to +\infty$)
\begin{equation}\label{eq:sol1}
	\eta(r)=2\sqrt{3} \arctan\left(e^{2(r-r_0)/L} \right)
	\ ,\qquad 
	e^{2A(r)}=\frac{e^{2a_0}}{\sin\left(\frac{\eta(r)}{\sqrt{3}}\right)}\ .
\end{equation}
The solution interpolates between two asymptotically $\text{AdS}_{4+1}$ spacetimes of the same radius $L$
\begin{equation}
	e^{2A} \underset{r\to-\infty}{\sim} \frac{e^{2a_0}}{2}e^{-2(r-r_0)/L}
	\ ,\qquad 
	e^{2A} \underset{r\to +\infty}{\sim} \frac{e^{2a_0}}{2}e^{2(r-r_0)/L}\ .
\end{equation}
The warp factor reaches its minimal value at $r=r_0$, $\eta(r_0)=\pi\sqrt{3}/2$, $e^{2A(r_0)}=e^{2a_0}$.

	\item The analytic extension $\phi\to \phi=i\, \eta$ makes the superpotential periodic in $\eta$, with period $4\pi$,
\begin{equation}
	\label{eq:Wex2}
	W(\eta)=\frac{4}{L}\cos\left(\frac{\eta}{2}\right)\ .
\end{equation}
There is a flow between the two critical points of the superpotential at $\eta=0$ (when $r\to-\infty)$ and $\eta=2\pi$ (when $r\to +\infty$)
\begin{equation}\label{eq:sol2}
	\eta(r)=4 \arctan\left(e^{(r-r_0)/L} \right)
	\ ,\qquad e^{2A(r)}=\frac{e^{2a_0}}{\sin^2\left(\frac{\eta(r)}{2}\right)}\ .
\end{equation}
The solution interpolates between two asymptotically $\text{AdS}_{3+1}$ spacetimes of the same radius $L$
\begin{equation}
	e^{2A} \underset{r\to-\infty}{\sim} \frac{e^{2a_0}}{4}e^{-2(r-r_0)/L}
	\ ,\qquad 
	e^{2A} \underset{r\to +\infty}{\sim} \frac{e^{2a_0}}{4}e^{2(r-r_0)/L}\ .
\end{equation}
The warp factor reaches its minimal value at $r=r_0$, $\eta(r_0)=\pi$, $e^{2A(r_0)}=e^{2a_0}$.

	\item The analytic extension $\phi^1\to i\eta^1                     $ renders the superpotential periodic in $\eta$, with period $2\sqrt{2}\pi$
\begin{equation}\label{eq:supex3}
	W(\eta^1,\phi^2) = \frac{2}{L}e^{-\frac{\phi^2}{\sqrt{6}}} \left[2+e^{\sqrt{\frac{3}{2}}\phi^2}\cos\left(\frac{\eta^1}{\sqrt{2}}\right)\right]\ .
\end{equation}
The solution can be found in implicit form, using $\eta^1$ as radial coordinate,
\begin{subequations}
\begin{eqnarray}
	e^{\sqrt{\frac{3}{2}}\phi^2(\eta^1)} &=&\cos\left(\frac{\eta^1}{\sqrt{2}}\right)+\sin^2\left(\frac{\eta^1}{\sqrt{2}}\right)\left\{C+\log\left[\text{cotan}\left(\frac{\eta^1}{2\sqrt{2}}\right)\right]\right\}
	\ ,\\
	e^{2A(\eta^1)} &=& e^{2a_0}\frac{e^{\sqrt{\frac{2}{3}}\phi^2(\eta^1)}}{\sin^2\left(\frac{\eta^1}{\sqrt{2}}\right)}\ .
\end{eqnarray}
\end{subequations}
The solution interpolates between $\eta^1=0$ (for $r\to -\infty$) and a finite value $0<\eta_0<\sqrt{2}\pi$ determined by the integration constant
\begin{equation}
	C=-\log\left[\text{cotan}\left(\frac{\eta_0}{2\sqrt{2}}\right)\right]-\frac{\cos\left(\frac{\eta_0}{\sqrt{2}}\right)}{\sin^2\left(\frac{\eta_0}{\sqrt{2}}\right)}\ .
\end{equation}
One has $\eta_0 \to 0$ for $C\to -\infty$ and $\eta_0 \to \sqrt{2}\pi$ for $C\to +\infty$.

This describes a flow from asymptotically $\text{AdS}_{4+1}$ space of radius $L$ and a singular geometry
\begin{equation}
	\eta \underset{r\to-\infty}{\sim} \eta_b e^{r/L}
	\ ,\qquad
	e^{2A} \underset{r\to-\infty}{\sim} \frac{2e^{2a_0}}{\eta_b^2}e^{-2r/L}\ ,
\end{equation}
and
\begin{align}
	&\dd s^2\sim \frac{L^2}{2\sin^2\left(\frac{\eta_0}{\sqrt{2}}\right)} e^{-2\sqrt{\frac{2}{3}}\phi}\dd\eta^2+\frac{e^{2a_0}}{\sin^2\left(\frac{\eta_0}{\sqrt{2}}\right)}e^{\sqrt{\frac{2}{3}}\phi}\eta_{\mu\nu}\dd x^\mu \dd x^\nu
	\ , \\ 
	&\phi\sim \sqrt{\frac{2}{3}}\log(\eta_0-\eta) \to -\infty\ .
\end{align}

	\item In this case the domain wall solution has only been found numerically. A preliminary examination indicates that the problem is equally complicated after the analytic extension, so there are no analytic solutions that can be found straightforwardly and we have not attempted to construct numerical solutions in this example.
    
	\item In this case there are several possibilities. We choose to do the analytic extension of both scalar fields $\{\phi^1,\phi^2\}\to \{i\eta^1,i\eta^2\}$. The superpotential is periodic, with periods $2\sqrt{3}\pi$ for $\eta^1$ and $2\pi$ for $\eta^2$:
\begin{equation}
	W(\eta^1,\eta^2) = \frac{3}{L}\left[\cos\left({\frac{\eta^1}{\sqrt{3}}}\right)+\cos\left({\eta^2}\right)\right]\ .
\end{equation}
There is a flow interpolating between two critical points of the superpotential, the origin $(\eta^1,\eta^2)=(0,0)$ (for $r\to -\infty$) and the point $(\eta^1,\eta^2)=(\sqrt{3}\pi,\pi)$ (for $r\to+\infty)$ 
\begin{subequations}\label{eq:sol5}
	\begin{equation}
	\eta^1(r) = 2\sqrt{3}\arctan{\left(e^{(r-r_1)/L}\right)}
	\ , \qquad
	\eta^2(r) = 2\arctan{\left(e^{3(r-r_2)/L}\right)}\ ,
\end{equation}
\begin{equation}
	e^{2A(r)}=e^{2a_0}\cosh\left(\frac{r-r_1}{L}\right)\left[\cosh\left(3~\frac{r-r_2}{L}\right)\right]^{1/3}\ .
\end{equation}
\end{subequations}
The solution interpolates between two asymptotically $\text{AdS}_{4+1}$ spacetimes of the same radius $L$
\begin{equation}
	e^{2A} \underset{r\to-\infty}{\sim} \frac{e^{2a_0}}{2^{4/3}}e^{-(2r-r_1-r_2)/L}
	\ ,\qquad 
	e^{2A} \underset{r\to +\infty}{\sim} \frac{e^{2a_0}}{2^{4/3}}e^{(2r-r_1-r_2)/L}\ .
\end{equation}    
	\item  The analytic extension $\phi\to i\eta$ renders the superpotential periodic in $\eta$, with period $2\pi$
\begin{equation}
	W(\eta)= \frac{6}{L}\cos^2\left(\frac{\eta}{2}\right)=\frac{3}{L}\left[1+\cos\left(\eta\right)\right]\ .
\end{equation}
There is a flow between the origin $\eta=0$ (at $r\to-\infty$) and the critical point of the superpotential at $\eta=\pi$ (at $r\to+\infty$)
\begin{equation}
	\eta(r) = 2\arctan{\left(e^{3 (r-r_0)/L}\right)}
	\ ,\qquad 
e^{2A(r)}=e^{2a_0} \left( \sin\frac{\eta(r)}{2}\right)^{-2/3}\ . 
\end{equation}
This flow interpolates between an asymptotically $\text{AdS}_{4+1}$ spacetime and a spacetime that is asymptotically flat
\begin{equation}
	e^{2A} \underset{r\to-\infty}{\sim} e^{2a_0}e^{-2(r-r_0)/L},\quad e^{2A} \underset{r\to +\infty}{\sim} e^{2a_0}\ .
\end{equation}
This solution is particularly interesting in the sense that the uplift to type IIB supergravity \cite{Gubser:2009qm,Cassani:2010uw,Gauntlett:2010vu} retains a real metric, while some of the fluxes supporting it are purely imaginary. It admits any Sasaki--Einstein manifold as internal geometry, including the five-sphere. The metric in ten-dimensions reads 
\begin{align}
\dd s^2_{10}=\cos\left(\frac{\eta}{2}\right)\dd s_5^2+\frac{L^2}{\cos\left(\frac{\eta}{2}\right)}\dd s_4^2\left(\text{\small KE}\right)+L^2\cos\left(\frac{\eta}{2}\right)\eta_{\text{\tiny SE}}^2\ ,
\end{align}
where $\dd s_5^2$ is the five-dimensional metric \eqref{eq:domain_wall}, $\dd s_4^2\left(\text{\small KE}\right)$ is the K\"ahler--Einstein base of the Sasaki--Einstein manifold and $\eta_{\text{\tiny SE}}$ describes the ${\rm U(1)}$ fiber over it. This metric is AdS when $\eta\to0$, but is singular in the limit $\eta\to \pi$ in which the five-dimensional metric is simply Minkowski. It is supported by both five and three-form fluxes, while the axio-dilaton vanishes. In terms of the standard forms defining the Sasaki--Einstein structure (see \cite{Gubser:2009qm,Cassani:2010uw,Gauntlett:2010vu} for details) the two-forms are
\begin{align}
B_2=\im\, L^2\tan\left(\frac{\eta}{2}\right){\rm Re}\,\Omega\,,\qquad\qquad C_2=-\im\,L^2\tan\left(\frac{\eta}{2}\right){\rm Im}\,\Omega\ ,
\end{align}
which are purely imaginary. The self-dual five-form flux, on the other hand, is real and given by
\begin{align}
F_5=L^4\left[2+3\tan^2\left(\frac{\eta}{2}\right)\right]\left(1+*\right)J\wedge J\wedge\eta_{\text{\tiny SE}}\ .
\end{align}

\item The analytic extension $\phi\to i\eta$ renders the superpotential periodic in $\eta$, with period $2\pi$
\begin{equation}
	W(\eta)= \frac{2}{L}\left[\cos(\eta)+1\right]\ .
\end{equation}
There is a flow between the origin $\eta=0$ (at $r\to-\infty$) and the critical point of the superpotential at $\eta=\pi$ (at $r\to+\infty$)
\begin{equation}
	\eta(r) = 2\arctan{\left(e^{2 (r-r_0)/L}\right)}
	\ , \qquad
e^{2A(r)}=e^{2a_0} \left( \sin\frac{\eta(r)}{2}\right)^{-1}\ . 
\end{equation}
This flow interpolates between an asymptotically $\text{AdS}_{3+1}$ spacetime and a spacetime that is asymptotically flat
\begin{equation}
	e^{2A} \underset{r\to-\infty}{\sim} e^{2a_0}e^{-2(r-r_0)/L}
	\ ,\qquad 
	e^{2A} \underset{r\to +\infty}{\sim} e^{2a_0}.
\end{equation}
The solution can be uplifted to eleven dimensions on any Sasaki--Einstein manifold using the results in \cite{Gauntlett:2009dn}. The metric reads
\begin{align}
\dd s^2=\left[\cos\left(\frac{\eta}{2}\right)\right]^{4/3}\dd s_4^2+\frac{4L^2}{\left[\cos\left(\frac{\eta}{2}\right)\right]^{2/3}}\dd s_6^2\left(\text{\small KE}\right)+4L^2\left[\cos\left(\frac{\eta}{2}\right)\right]^{4/3}\eta_{\text{\tiny SE}}^2\ ,
\end{align}
where $\dd s_4^2$ is the domain-wall metric  \eqref{eq:domain_wall}, $\dd s_6^2\left(\text{\small KE}\right)$ is the metric of a six-dimensional K\"ahler--Einstein base and $\eta_{\text{\tiny SE}}$ is the U(1) fiber over it. There is an IR singularity as $\eta\to\pi$. This metric is supported by the four-form flux  
\begin{align}
F_4&=\frac{3}{L}\left[1+\frac43\tan^2\left(\frac{\eta}{2}\right)\right]\cos^4\left(\frac{\eta}{2}\right){\rm vol}_4-32\im L^3\tan\left(\frac{\eta}{2}\right){\rm Re}\,\Omega\wedge\eta_{\text{\tiny SE}}\nonumber\\[2mm]
&+16\im L^2\tan\left(\frac{\eta}{2}\right)\dd r\wedge{\rm Im}\,\Omega\ ,
\end{align}
with ${\rm vol}_4$ the volume form of the four-dimensional metric \eqref{eq:domain_wall}. The definition and properties of the one-form $\eta_{\text{\tiny SE}}$ and complex three-form $\Omega$ can be found in \cite{Gauntlett:2009dn}. Notice that the internal part is purely imaginary, while the spacetime components are real. 

\item Performing the analytic extension $\phi^1\to\im\eta^1$ while keeping $\phi^2$ real results in a periodic superpotential with period $2\sqrt{2}\pi$
\begin{align}
W(\eta^1,\phi^2)=\frac{2}{L}e^{-\frac{\phi^2}{\sqrt{6}}}\left[e^{\sqrt{\frac{3}{2}}\phi^2}+2\cos\left(\frac{\eta^1}{\sqrt{2}}\right)\right]\,.
\end{align}
The structure of this superpotential is very similar to that in \eqref{eq:supex3}. The solution can be given again using $\eta^1$ as radial coordinate
\begin{align}
	e^{2A(\eta^1)}=e^{2a_0}\frac{e^{-\frac{\phi^2(\eta^1)}{\sqrt{6}}}}{\sin\left(\frac{\eta^1}{\sqrt{2}}\right)}
	\ ,\qquad e^{\sqrt{\frac32}\phi^2(\eta^1)}=\frac{1}{\cos\left(\frac{\eta^1}{\sqrt{2}}\right)+C\sin\left(\frac{\eta^1}{\sqrt{2}}\right)}\,.
\end{align}
The radial coordinate ranges between $\eta^1=0$ and $\eta^1=\eta_0$, where the upper value depends on the integration constant as $C=-\cot\left(\eta_0/2\right)$. It interpolates between ${\rm AdS}_{3+1}$ and a hyperscaling-violating geometry with coefficient $\theta=4$.

\end{enumerate}

The solutions in examples 1, 2 and 5, interpolating between two asymptotically AdS spaces, can be thought of as instances of ``traversable wormholes'' where the only way to travel between the asymptotic regions is through the wormhole itself. Solutions 1 and 2 turn out to be particular examples of the families studied in \cite{Huang:2020qmn,Nozawa:2020gzz}.

Without a violation of the null energy condition, this type of solutions should not be possible \cite{Morris:1988tu,Hochberg:1998ii}. Nonetheless, such violation does not lead to causality violations of the type discussed in \cite{Morris:1988tu}, since it is not possible to travel between the two asymptotic regions outside the wormhole. There are also no causality violations for signals connecting two points at the boundary, of the type discussed in \cite{Hoyos:2010at}, as the local speed of light is the same in all the geometry, this latter being manifest from the domain-wall ansatz \eqref{eq:domain_wall}. In the rest of the paper we will focus on these solutions, as they are the most interesting for their connection to the holographic description of entanglement and to the quantum gravitational path integral.

\subsection{Perturbative stability and fluctuation spectrum of wormholes}

Even though the solutions we are constructing satisfy a set of BPS equations and can be seen as analytic extensions of supersymmetric solutions, their stability is not guaranteed due to the change in sign of the kinetic term of the analytically continued scalars. In this subsection, we study the perturbative stability of backgrounds with only ghost scalars coupled to gravity. More complicated cases with both ghost and ordinary scalars are left for future studies.

We expand around the background domain wall solution to linear order in the metric and scalar fluctuations. We enforce the radial gauge and consider a plane wave expansion. We then identify the combination of fluctuating fields that are invariant under the remaining diffeomorphisms that are not fixed by the gauge choice. Full details can be found in Appendix \ref{app:eoms}. For invariant fluctuations involving only the metric components, the equations of motion take the form derived in \eqref{eq:invariantHeq}, in $d+1$ dimensions,
\begin{equation}
	H''+(d-4)A' H'+\left[e^{-2A}M^2-2(d-2)(A')^2-\frac{1}{d-1}(\eta_0')^2\right]H=0\ ,
\end{equation}
with $H$ the invariant fluctuation and $M^2=-\eta^{\mu\nu}k_\mu k_\nu$ the mass for a mode with momentum $k_\mu$ along the field theory directions. The radial functions $\eta_0(r)$ and $A(r)$ are the background values of the ghost scalar and the warp factor, respectively. 

It is convenient to write the equation in Sturm-Liouville form
\begin{equation}
	\frac{d}{dr}\Big[p(r)H'(r)\Big]+q(r)H(r)+M^2 w(r) H(r)=0\ ,
\end{equation}
with
\begin{equation}
	p(r)=e^{(d-4)A}
	,\quad 
	q(r)=-e^{(d-4)A}\left[2(d-2)(A')^2+\frac{1}{d-1}(\eta_0')^2\right],\quad w(r)=e^{(d-6)A}.
\end{equation}
Normal modes are solutions to the equations that are normalizable at the aymptotic boundaries (in the case of wormholes), or normalizable at the asymptotic boundary and regular in the interior. This is a Sturm-Liouville problem with eigenvalues $\lambda=M^2$. Non-trivial solutions will typically exist only for a discrete set of values $\lambda_n=M_n^2$, $n=1,2,3,\dots$. Since these are real equations with real coefficients and the boundary conditions are also real, the modes and eigenvalues will also be real. Therefore, an instability will happen iff there is an eigenmode with a negative eigenvalue $\lambda_n=M_n^2<0$. The modes in the spectrum split according to parity in the holographic radial direction $r\to -r$, if the geometry has this symmetry.

Now we show that all eigenvalues are positive. For this, we multiply the equation in Sturm-Liouville form by the eigenfunction and integrate over the radial coordinate. We find, for any $n$,
\begin{equation}
	\int_a^b  \dd r \, \left[H_n\frac{d}{dr}\Big(p H_n'\Big)+q H_n^2+M_n^2 w H_n^2\right]=0\ ,
\end{equation}
where the integration domain $(a,b)$ spans the entire range of the radial coordinate. We can integrate by parts the first term inside the integral. For a normal mode, the resulting boundary term should vanish at the extremes of the interval, so after moving some terms to the other side of the equality, we are left with
\begin{equation}
	M_n^2\int_a^b  \dd r \, w H_n^2=\int_a^b  \dd r \, \left[p (H_n')^2-q H_n^2\right]\ .
\end{equation}
Since $H_n$ is real, and $p>0$, $q<0$, and $w>0$ in the full range of the integration, the integrals at both sides of the equality are positive definite and it follows that $M_n^2>0$.
This proves that these backgrounds are perturbatively stable as far as the $H$ fluctuations are concerned. Curiously, this argument does not work for an ordinary scalar, since the term in the coefficient $q$ proportional to $(\eta_0')^2$ would have the opposite sign, and $q$ would not have a definite sign.

To complete the stability analysis, we need to study the remaining gauge-invariant fluctuation $Z_\eta$, the one which involves the ghost scalar. We define it in \eqref{eq:zetaeta} and its equation of motion is derived in \eqref{eq:invariantZeq},
\begin{equation}
	\label{eq:eom_Zeta}
	0=Z_\eta''+dA' Z_\eta'+\left(e^{-2A}M^2+Q_{ZZ}\right)Z_\eta+Q_{ZH}(H_S'-2A' H_S).
\end{equation}
We remind the reader that $H_S$ is a gauge-invariant, scalar mode that involves only the metric fluctuations \eqref{eq:HS}. 

One can study the spectrum of \eqref{eq:eom_Zeta} setting $H_S=0$. In fact, $Z_\eta$ does not source other modes and we have already determined that the spectrum of $H_S$ is stable. Unfortunately, the coefficient $Q_{ZZ}$ given in \eqref{eq:Qcoefs} and \eqref{eq:Qcoef2}, does not have a definite sign, so the Sturm-Liouville argument that we used above does not apply and it is necessary to compute the spectrum of $Z_\eta$ explicitly. We have performed the calculation of the first eigenvalues in the spectrum for examples 1, 2, and 5\footnote{Note that example 5 features two scalar ghosts, but this does not produce qualitative changes in the analysis with respect to the cases with just a single scalar ghost.} from Subsection \ref{subsec:examples}. For convenience, we have adopted different approaches, Chebyshev spectral method, analytic analysis, and numerical shooting, respectively. In the three cases that we examined, we found no evidence of instability (see Appendix \ref{app:numerics} for further details).

The scalar fluctuations present some interesting qualitative differences with respect to the metric fluctuations. In fact, the ghost scalar modes split in two independent sectors with the same masses, each of them localized at one side of the wormhole. For all regular modes, the value of the field and its derivative vanishes at the center of the wormhole, so the fluctuation can be consistently set to zero on the other side. As a consequence, the two-point correlation function mixing the two copies of the operator dual to the ghost field will vanish with ordinary boundary conditions.

